\documentclass[11pt]{article}
\usepackage{amsmath,amssymb,amsthm,hyperref}
\newtheorem{theorem}{Theorem}
\newtheorem{lemma}{Lemma}
\newtheorem{proposition}{Proposition}
\newtheorem{corollary}{Corollary}
\newcommand{\Prb}{\mathbb P}
\newcommand{\E}{\mathbb E}
\newcommand{\one}{\mathbf 1}
\title{Small row supports and identification of an unknown equitable frame}
\author{Independent support audit}
\date{15 November 2026}
\begin{document}
\maketitle

\paragraph{Scope.}
This note proves the small-support and deterministic compatibility steps
of the root's candidate coarse-frame recovery argument. It does not audit
the entropy/type enumeration or prove its required macroscopic exclusion
statement. In particular it does not claim the full proposed recovery
theorem. All counts below allow arbitrary support sizes and compositions
across planted blocks; no balanced-support assumption is made.

\paragraph{Model.}
Partition $n$ labeled vertices into fixed blocks $B_1,\ldots,B_k$, where
$|B_i|=N_i\geq\rho n$ for a fixed $\rho>0$. A fixed nonnegative integer
degree frame $Q=(d_{ij})$ specifies $d_{ij}$ stubs from each vertex in
$B_i$ toward block $B_j$, with $N_i d_{ij}=N_j d_{ji}$ for $i\ne j$.
Internal stubs are paired uniformly, independently for different blocks;
for $i<j$, their two equal cross-stub pools are matched by an independent
uniform bijection. Assume all parity conditions needed to define these
factors. Zero-degree factors are omitted. Write
\[
 \delta=\min_i\sum_jd_{ij}\geq3,
 \qquad D=\max_i\sum_jd_{ij}.
\]
The configuration graph may have loops or multiple edges. A loop counts
twice in degrees and contributes two to the adjacency diagonal. This
gives the exact identity $AJ=JQ$ for the planted indicator matrix $J$.
Internal degrees at least two are not needed for the results of this note;
they are an additional assumption in the root's separate entropy argument.

Call $S\subseteq V$ \emph{bad} if no vertex outside $S$ has exactly one
incident edge into $S$, counting parallel edges with multiplicity.

\begin{theorem}[Uniform small-set bound in the typed configuration model]
Put $c=(\delta-2)/2$ and
\[
 C=4(D+1)^2\exp(1+D/2+1/e)\,
       2^{D^2/2}(4D/\rho)^D,
\]
and choose
\[
 0<\eta\leq
 \min\left\{\frac{\rho}{4D},\ (2Ce^c)^{-1/c}\right\}.
 \tag{1}
\]
Let $Z_s$ be the number of bad subsets of size $s$. For every integer
$1\leq s\leq\eta n$,
\[
 \E Z_s\leq C^s\left(\frac{s}{n}\right)^{\lceil cs\rceil}.
 \tag{2}
\]
Consequently,
\[
 \begin{split}
 \Prb\{\exists\hbox{ bad }S:1\leq|S|\leq\eta n\}
 &\leq\sum_{s=1}^{\lfloor\eta n\rfloor}
             C^s(s/n)^{\lceil cs\rceil}\\
 &\leq Cn^{-\lceil c\rceil}
       +2C^2(2/n)^{\delta-2}
 =O_{D,\rho}(n^{-\lceil(\delta-2)/2\rceil}).
 \end{split}
 \tag{3}
\]
The constants are uniform in the distribution of $S$ over the planted
blocks, and over any frames satisfying the stated degree and block-size
bounds. The total number of frame types does not enter the estimate.
\end{theorem}
\begin{proof}
Fix a bad set $S$ of size $s$, and let $W$ be the set of its outside
neighbors, of size $w$. Let $a$ count edges inside $S$, including each
loop as one edge, and let $b$ count edges between $S$ and $W$. Put
$e=a+b$. The sum of degrees in $S$ is $2a+b$, so
\[
 \delta s\leq2a+b\leq Ds.
\]
Every vertex of $W$ has at least two edges into $S$, giving $2w\leq b$.
In particular $w\leq Ds/2$, $e\leq Ds$, and
\[
 e-s-w=a+b-s-w
 \geq a+b/2-s\geq(\delta-2)s/2=cs.
 \tag{4}
\]
The left side is an integer, hence at least $\lceil cs\rceil$.

A witness specifies $S,W$, all stubs of $S$, selected stubs of $W$, and
their $e$ prescribed stub pairs. Every positive typed factor has at least
$\rho n$ stubs in each relevant pool. If a witness prescribes $m$ pairs
in an internal pool of size $L$, its probability is
\[
 \frac1{(L-1)(L-3)\cdots(L-2m+1)}.
\]
For a cross-factor bijection the probability is $1/(L)_m$.
As $m\leq e\leq Ds$ and $s\leq\rho n/(4D)$, each denominator factor
is at least $\rho n/2$. Independence of typed factors therefore bounds
the probability of any valid prescribed partial configuration by
\[
 (2/(\rho n))^e.                                  \tag{5}
\]
Incompatible typed prescriptions have probability zero and can be counted
as well for an upper bound. This also covers all loop prescriptions.

There are at most $\binom ns\binom nw$ choices of $S,W$ (overcounting
intersections is harmless). The selected stubs at $W$ have at most
$2^{Dw}$ possibilities. The $2e$ selected stubs can be paired in at most
$(2e-1)!!\leq(2Ds)^e$ ways. We may count pairings not respecting the
types or the requirement that each edge touch $S$; true witnesses form a
subset. Combining this count with (5), for each permitted $e,w$ the
expectation contribution is at most
\[
 \binom ns\binom nw\,2^{Dw}(2Ds)^e(2/(\rho n))^e.
 \tag{6}
\]
For $w>0$, use the usual binomial bounds and rewrite (6) as at most
\[
 e^{s+w}2^{Dw}(4D/\rho)^e
 (s/n)^{e-s-w}(s/w)^w.
\]
For $w=0$ the same expression is read with $(s/w)^w=1$.
The elementary maximum of $t\log(1/t)$ is $1/e$, so
$(s/w)^w\leq\exp(s/e)$. Also $w\leq Ds/2$, $e\leq Ds$, and (4) holds.
Thus the contribution is bounded by
\[
 \left[\exp(1+D/2+1/e)\,2^{D^2/2}(4D/\rho)^D\right]^s
 (s/n)^{\lceil cs\rceil}.
\]
There are at most $(Ds+1)^2\leq[4(D+1)^2]^s$ possible pairs $e,w$.
This proves (2), uniformly without any restrictions on the types of $S$.

For the sum, retain the integer exponent for $s=1$. For $s\geq2$ put
$u_s=C^s(s/n)^{cs}$, which is an upper bound for the summand. Whenever
$s+1\leq\eta n$,
\[
 \frac{u_{s+1}}{u_s}
 =C\left(\frac{s+1}{n}\right)^c
      (1+1/s)^{cs}
 \leq Ce^c\eta^c\leq\frac12.
\]
Hence $\sum_{s=2}^{\lfloor\eta n\rfloor}u_s\leq2u_2$, with an empty
sum interpreted as zero. Adding the $s=1$ term yields (3). Since
$\delta-2\geq\lceil(\delta-2)/2\rceil$, the final order bound follows.
\end{proof}

\begin{corollary}[Uniform simple factors]
Condition all typed factors to be simple. Their joint simplicity
probability is bounded below by a positive constant $p_0$ for fixed
degrees and fixed positive block proportions. Then the resulting factors
are independent uniform simple regular or biregular factors, and
\[
 \Prb_{\rm simple}\{\exists\hbox{ bad }S:
                    1\leq|S|\leq\eta n\}
 \leq p_0^{-1}2C^2(2/n)^{\delta-2}=O(n^{-(\delta-2)}).
 \tag{7}
\]
\end{corollary}
\begin{proof}
Every simple factor has the same number of labeled-stub realizations.
Thus conditioning gives precisely the claimed uniform simple model.
The positive simplicity probability is the usual fixed-degree
configuration-model fact, applied to finitely many independent factors.
For reference, the limiting probabilities are
$\exp[-(a^2-1)/4]$ for an internal $a$-regular factor and
$\exp[-(u-1)(v-1)/2]$ for a cross $(u,v)$-biregular factor; their product
is positive. Only a positive lower bound is needed here.

A singleton set cannot be bad in a loopless simple graph of positive
degree: every neighbor outside it has exactly one edge into it. Thus
the $s=1$ term disappears on the simplicity event. Bound the intersection
of that event with the union over $s\geq2$ by the corresponding
unconditioned witness sum, and divide by its probability at least $p_0$.
The geometric-series bound in the theorem proves (7).
\end{proof}

\begin{corollary}[Uniform exclusion of small row supports]
On the event that no bad set of size at most $\eta n$ exists, there is
no nonzero matrix $V$ over any field with at most $\eta n$ nonzero rows
satisfying $AV=VR$, for any square matrix $R$ of any finite order.
Thus the exclusion is simultaneous in $R$, the number of columns, and
the entries of $V$.
\end{corollary}
\begin{proof}
Let $S$ be the nonzero row support. At a vertex $v$ outside $S$, the
equation gives $\sum_uA_{vu}V_u=0$. If exactly one edge joined $v$ to
$S$, this sum would be one nonzero row $V_u$, a contradiction. Loops
outside $S$ multiply the zero row and do not matter. Therefore $S$ would
be bad, contradicting the stated event. No assumption on $R$ is used.
\end{proof}

\begin{proposition}[Exact near-point compatibility]
Let $Y\in\{0,1\}^{n\times r}$ have exactly one nonzero entry in every
row and satisfy the INTEGER equation $AY=YR$. Let
$T\in\{0,1\}^{k\times r}$ also have one nonzero entry per row. Define
$S=\{v:Y_v\ne(JT)_v\}$, $s_i=|S\cap B_i|$. If
\[
 N_i>s_i+\sum_jd_{ji}s_j\quad\hbox{for every }i,    \tag{8}
\]
then $QT=TR$. In particular, it is enough that
\[
 |S|<\frac{\min_iN_i}{D+1}.
 \tag{9}
\]
Under either condition $V=Y-JT$ satisfies $AV=VR$.
\end{proposition}
\begin{proof}
In block $B_i$, at most $s_i$ vertices are changed themselves. The
number of vertices adjacent to a changed vertex is at most
$\sum_jd_{ji}s_j$, since each changed vertex in block $j$ has exactly
$d_{ji}$ incident edges toward block $i$. This is an upper bound even
with loops or parallel edges. Condition (8) supplies a vertex $v\in B_i$
whose own row and all neighbor rows are unchanged. At that vertex,
\[
 (QT)_i=(AJT)_v=(AY)_v=(YR)_v=(TR)_i.
\]
This holds for every $i$. The global bound
$s_i+\sum_jd_{ji}s_j\leq(D+1)|S|$ proves (9) suffices. Finally
\[
 A(Y-JT)=YR-JQT=YR-JTR=(Y-JT)R.
\]
\end{proof}

Choose the small-support constant additionally to satisfy
$\eta\leq\rho/[2(D+1)]$. Combining this proposition with the uniform
small-support event shows that every feasible $Y$ within $\eta n$ row
changes of a planted point map $JT$ is EXACTLY that point map. This
argument requires neither balanced error sizes in the planted blocks nor
knowledge of the candidate quotient matrix. For a merely modular equation,
the clean-vertex step only proves $QT=TR$ in that field; the integer
conclusion applies to the candidates verified over the integers.

\begin{proposition}[No tiny classes in a bounded connected quotient]
Let $Y$ encode a partition into $r$ nonempty classes of sizes
$M_1,\ldots,M_r$, and suppose $AY=YR$ for a nonnegative integer matrix
$R$ whose row sums are at most $D\geq1$. Suppose the off-diagonal support
of $R$ is connected after forgetting orientations. Then
\[
 \min_aM_a\geq\frac{n}{rD^{r-1}}
 \geq\frac{n}{KD^{K-1}}\qquad(r\leq K).
 \tag{10}
\]
\end{proposition}
\begin{proof}
Undirected edge counting gives $M_aR_{ab}=M_bR_{ba}$ for $a\ne b$.
Since all classes are nonempty, $R_{ab}>0$ if and only if $R_{ba}>0$.
Along a path of length $\ell\leq r-1$, successive class-size ratios
are ratios of positive integers at most $D$, hence between $D^{-1}$
and $D$. Starting from a smallest class, every class has size at most
$D^{r-1}$ times it. Sum the sizes to obtain (10). Diagonal entries,
including contributions from loops, do not affect this argument.
\end{proof}

\begin{proposition}[Identification of unknown $k$ and frame, conditional
on the macroscopic step]
Suppose the true planted frame $Q$ is connected, has $k\leq K$ nonempty
blocks, and has row sums at most the fixed bound $D$. Consider every
nonnegative integer $r\times r$ frame $R$ with $r\leq K$, row sums at
most $D$, and connected off-diagonal support, and every verified nonempty
one-hot solution $Y$ of $AY=YR$. Assume the separate macroscopic argument
proves that with probability $1-o(1)$ every such $Y$ lies within
$\eta n$ changed rows of SOME planted point map $JT$.

Then, with probability $1-o(1)$, all accepted partitions are equitable
coarsenings of the planted partition. Selecting any accepted partition
with the largest number of classes recovers $k$ and the planted partition
up to class permutation. Its frame is then recovered by exact neighbor
counts. This implication uses no entropy assertion beyond the stated
macroscopic hypothesis.
\end{proposition}
\begin{proof}
Intersect the macroscopic event with the single uniform small-support
event, and apply near-point compatibility followed by small-support
exclusion. It gives $Y=JT$ for every accepted candidate simultaneously.
Nonempty candidate classes force $T$ to be onto, hence $r\leq k$.
The planted pair $(J,Q)$ is among the candidates, so the maximum is $k$.
At $r=k$, an onto map of $k$ old types to $k$ labels is a bijection,
which recovers their blocks up to label permutation. Given the blocks,
each row of $R$ is their integer neighbor-count vector and is unique.

There are at most $\sum_{r=1}^K(D+1)^{r^2}$ possible frame matrices.
Thus a union over candidate frames is finite with constants independent
of $n$. The size bound (10) ensures no vanishing nonempty class escapes
a compactness argument that needs it. The uniform small-support event
itself does not require any union bound over candidate frames.
\end{proof}

\paragraph{Limits and necessary qualifications.}
The minimum total degree of three is essential for the present small-set
conclusion: fixed-degree two configuration graphs have small cycle
components, which are bad sets. Connectedness of a candidate quotient
and integrality of its entries are essential to the class-size path
argument; bounded real positive entries would not bound reverse ratios.
The unknown-frame procedure requires a fixed finite upper bound $K$.
No claim about efficiently determining arbitrary unbounded numbers of
classes follows here. No enumeration-runtime, entropy-deficit, or priority
claim is supplied by this note.
\end{document}
